% !TeX root = resume.tex
% Condensed from the existing project entry; retain methods and measured results.

\projectheading{小模型 Tool Use / Function Calling 能力提升研究}{2026.05\,--\,至今}
\textbf{项目介绍：}针对小模型工具误选、参数幻觉与无关工具误调用，基于 BFCL v4 开展数据蒸馏、SFT 消融与 GRPO 优化实验。\par
\begin{experienceitems}
  \item \textbf{数据构建：}基于 xLAM-function-calling-60k 构建 8 种工具描述／输出格式变体，以 DeepSeek-V4、MiniMax-m2.5 蒸馏推理链与调用结果，经重试清洗形成各 2k／5k／10k 的 SFT 数据集。
  \item \textbf{SFT 消融：}在 Llama3.2-1B-Instruct 上完成两类教师、三档数据量及 prompt-only 基线共 7 组对照；MiniMax-5k 将 Non-Live Acc 从 38.4\% 提升至 \textbf{68.6\%（+30.2pp）}，Live Acc 从 11.4\% 提升至 \textbf{51.0\%（+39.6pp）}。
  \item \textbf{误调用治理：}针对正调用样本偏置，合成覆盖无关工具、参数缺失、参数冲突与约束不支持的 300 条 no-call 负样本。在 MiniMax-2k 基础上重训，Overall Irrelevance 从 3.48\% 提升至 \textbf{53.58\%（+50.1pp）}；同时记录 Live Acc 从 46.93\% 降至 40.04\%，评估拒用能力增强对正常调用的影响。
  \item \textbf{GRPO 对照：}在 Qwen2.5-3B-Instruct 上基于 TRL 与 vLLM，设计覆盖格式、工具名、参数名及参数值的细粒度奖励，完成 GRPO-only、SFT+GRPO、两组 SFT 及基线共 5 组实验；当前 GRPO 方案未超过基线，SFT 组基本持平。
\end{experienceitems}
